%% kapa: exam=1 round=32 year=2021 period=2 form=A subject=accounting subject_ko=회계학 seq=31 number=71
%% source: Q-Net 감정평가사 기출문제 (cst00309) artlSeq=5208397 file=extracted/1008/5208397/A형/2021년도 제32회 감정평가사 1차 2교시 A형.pdf
%% conversion: kordoc 4.24.51 -> tools/gen1.py (auto); review=proofread-sonnet; flags=none
\begin{question}{71}
원가관리기법에 관한 설명으로 옳은 것은?
\begin{choices}
  \item 제약이론을 원가관리에 적용한 재료처리량공헌이익(throughput contribution)은 매출액에서 기본원가를 차감하여 계산한다.
  \item 수명주기원가계산에서는 공장자동화가 이루어지면서 제조이전단계보다는 제조단계에서의 원가절감 여지가 매우 높아졌다고 본다.
  \item 목표원가계산은 표준원가와 마찬가지로 제조과정에서의 원가절감을 강조한다.
  \item 균형성과표는 전략의 구체화와 의사소통에 초점이 맞춰진 제도이다.
  \item 품질원가계산에서는 내부실패원가와 외부실패원가를 통제원가라 하며, 예방 및 평가활동을 통해 이를 절감할 수 있다.
\end{choices}
\end{question}
\answer{4}
